% Part: many-valued-logic
% Chapter: sequent-calculus
% Section: introduction

\documentclass[../../../include/open-logic-section]{subfiles}

\begin{document}

\olfileid{mvl}{seq}{int}

\olsection{پېژندنه}

د کلاسيک منطق د سېکوېنټ حساب د !!{derivation} يو اغېزمن او ساده نظام
دے۔ که ډېرارزښته منطق په داسې ماتريس تعريف شوے وي چې د رښتياوالي
قيمتونه ئې متناهي وي، يعنې~$V$ متناهي وي، د هغه دپاره سېکوېنټ حساب
ورکول کېدے شي۔ د دې کار مفکوره د کلاسيکې قضيې د سېکوېنټونو له معناوو
او د استنتاج د قاعدو له بڼې څخه راځي۔

اوس راياد کړئ چې يو سېکوېنټ
\begin{align*}
    !A_1, \dots, !A_m & \Sequent !B_1, \dots, !B_n
\intertext{د لاندې !!{formula} په توګه تعبيرېدے شي:}
    (!A_1 \land \cdots \land !A_m) & \lif (!B_1 \lor \cdots \lor
!B_n)
\end{align*}
د سرچينې يادونه: په اصلي متن کښې د سېکوېنټ د کيڼ لوري وروستے
شاخص n چاپ شوے، خو د همدغه تعبير د کيڼ اړخ وروستے شاخص m دے؛
دلته دواړه کيڼ شاخصونه يو شان کړل شول۔ په نورو ټکو کښې، !!^a{valuation}~$\pAssign{v}$
يو سېکوېنټ~$\Gamma \Sequent \Delta$ هله او يوازې هله \emph{پوره کوي} چې
يا د کوم~$!A \in \Gamma$ دپاره~$\pValue{v}(!A) = \False$
وي، يا د کوم~$!A \in \Delta$ دپاره~$\pValue{v}(!A) = \True$
وي۔ د دې تعبير له مخې ابتدايي سېکوېنټونه~$!A \Sequent !A$ تل پوره
کېږي، ځکه چې يا~$\pValue v(!A) = \True$ وي يا~$\pValue v(!A) =
\False$ وي۔ د سرچينې يادونه: د دې وروستۍ مساوات په اصلي متن کښې
د قيمت تابعې له نښې سره v نۀ دے راغلے؛ دلته څرګند شوے دے۔

دا په~$\Log{LK}$ کښې د شرطي نښلوونکي د استنتاج قاعدې دي، چې د څنګ
فارمولونه~$\Gamma$ او~$\Delta$ ترې اېستل شوي دي:

\begin{defish}
    \Axiom$ \fCenter !A$
    \Axiom$ !B \fCenter $
    \RightLabel{\LeftR{\lif}}
    \BinaryInf$ !A \lif !B \fCenter $
    \DisplayProof
    \hfill
    \Axiom$ !A \fCenter !B$
    \RightLabel{\RightR{\lif}}
    \UnaryInf$ \fCenter !A \lif !B $
    \DisplayProof
\end{defish}

که د سېکوېنټ پورته معنايي تعبير وکاروو، د~$\LeftR{\lif}$ قاعده
داسې لوستلے شو: که~$\pValue v(!A) = \True$ او~$\pValue v(!B) =
\False$ وي، نو~$\pValue v(!A \lif !B) = \False$ وي۔ همداسې،
د~$\RightR{\lif}$ قاعده وايي چې که يا~$\pValue v(!A) = \False$
وي يا~$\pValue v(!B) = \True$ وي، نو~$\pValue v(!A \lif !B) =
\True$ وي۔ په حقيقت کښې دا شرطي بيانونه دوه‌لوريز هم دي۔ د
$\LeftR{\land}$ او~$\RightR{\lor}$ د قاعدو په معياري بڼه کښې اړوند
شرطي بيانونه دوه‌لوريز نۀ وي۔ خو د دغو قاعدو بديلې بڼې شته چې
دوه‌لوريز دي:

\begin{defish}
    \Axiom$!A, !B, \Gamma \fCenter \Delta$
    \RightLabel{\LeftR{\land}}
    \UnaryInf$!A \land !B, \Gamma \fCenter \Delta$
    \DisplayProof
    \hfill
    \Axiom$ \Gamma \fCenter \Delta, !A, !B$
    \RightLabel{\RightR{\lor}}
    \UnaryInf$ \Gamma \fCenter \Delta, !A \lor !B$
    \DisplayProof
\end{defish}

دا بنسټيزه مفکوره چې پر يو~$n$ ارزښته منطق پلې شي، بيا داسې
سېکوېنټ حساب ورکوي چې د دوو پر ځای~$n$ ځايونه لري، د رښتياوالي د
هر قيمت دپاره يو ځاے۔ د درې‌ارزښته منطق دپاره، چې~$V = \{\False,
\Undef, \True\}$ لري، سېکوېنټ د~$\Gamma \mid \Pi \mid \Delta$
په بڼه عبارت دے۔ دا په~!!a{valuation}~$\pAssign v$ کښې هله پوره کېږي
چې يا د کوم~$!A \in \Gamma$ دپاره~$\pValue{v}(!A) =
\False$ وي، يا د کوم~$!A \in \Delta$ دپاره~$\pValue{v}(!A)
= \True$ وي، يا د کوم~$!A \in \Pi$ دپاره~$\pValue{v}(!A) =
\Undef$ وي۔ نو ابتدايي سېکوېنټونه~$!A \mid !A \mid !A$ تل پوره
کېږي۔

\end{document}
